% =========================================================
% 10_dot_cross_products.tex
% Vectors and Linear Algebra
% =========================================================

% =========================================================
% SECTION DIVIDER
% =========================================================

\section{Dot and Cross Products}

% =========================================================
% DOT PRODUCT
% =========================================================

\begin{frame}{Dot Product}

For two vectors
\vspace{-0.1cm}
\[
\vect{u}=
\begin{bmatrix}
u_1\\
u_2\\
\vdots\\
u_n
\end{bmatrix},
\qquad
\vect{v}=
\begin{bmatrix}
v_1\\
v_2\\
\vdots\\
v_n
\end{bmatrix},
\]

their \textbf{dot product} is
\vspace{-0.1cm}
\[
\boxed{
\vect{u}\cdot\vect{v}
=
\sum_{i=1}^{n}u_iv_i
}
\]

or equivalently,
\vspace{-0.1cm}
\[
\boxed{
\vect{u}\cdot\vect{v}
=
\vect{u}^{T}\vect{v}.
}
\]

The result is a \textbf{scalar}, not a vector.

\end{frame}


% =========================================================
% DOT PRODUCT EXAMPLE
% =========================================================

\begin{frame}{Example: Dot Product}

Let
\vspace{-0.1cm}
\[
\vect{u}
=
\begin{bmatrix}
2\\
-1\\
3
\end{bmatrix},
\qquad
\vect{v}
=
\begin{bmatrix}
1\\
4\\
2
\end{bmatrix}.
\]

Then
\vspace{-0.1cm}
\[
\vect{u}\cdot\vect{v}
=
(2)(1)+(-1)(4)+(3)(2).
\]

Therefore,

\[
\boxed{
\vect{u}\cdot\vect{v}=4.
}
\]

Using matrix notation,
\vspace{-0.1cm}
\[
\vect{u}^{T}\vect{v}
=
\begin{bmatrix}
2&-1&3
\end{bmatrix}
\begin{bmatrix}
1\\4\\2
\end{bmatrix}
=4.
\]

\end{frame}


% =========================================================
% GEOMETRIC DOT PRODUCT
% =========================================================

\begin{frame}{Geometric Meaning of the Dot Product}

For non-zero vectors,

\[
\boxed{
\vect{u}\cdot\vect{v}
=
\|\vect{u}\|
\|\vect{v}\|
\cos\theta
}
\]

where $\theta$ is the angle between the vectors.

Therefore,

\[
\boxed{
\cos\theta
=
\frac{\vect{u}\cdot\vect{v}}
{\|\vect{u}\|\|\vect{v}\|}
}
\]

The dot product therefore measures the directional
relationship between two vectors.

\end{frame}


% =========================================================
% ORTHOGONALITY
% =========================================================

\begin{frame}{Orthogonality}

Two vectors are \textbf{orthogonal} if they are perpendicular.

For non-zero vectors,
\vspace{-0.1cm}
\[
\vect{u}\perp\vect{v}
\]

precisely when
\vspace{-0.1cm}
\[
\boxed{
\vect{u}\cdot\vect{v}=0.
}
\]

For example,
\vspace{-0.1cm}
\[
\vect{u}
=
\begin{bmatrix}
1\\
2
\end{bmatrix},
\qquad
\vect{v}
=
\begin{bmatrix}
2\\
-1
\end{bmatrix}.
\]

Then
\vspace{-0.1cm}
\[
\vect{u}\cdot\vect{v}
=
(1)(2)+(2)(-1)
=
0.
\]

Hence,
\vspace{-0.1cm}
\[
\boxed{
\vect{u}\perp\vect{v}.
}
\]

\end{frame}


% =========================================================
% VECTOR PROJECTION
% =========================================================

\begin{frame}{Projection of One Vector onto Another}

The projection of $\vect{u}$ onto a non-zero vector $\vect{v}$ is

\[
\boxed{
\operatorname{proj}_{\vect{v}}\vect{u}
=
\frac{\vect{u}\cdot\vect{v}}
{\vect{v}\cdot\vect{v}}
\vect{v}.
}
\]

Since

\[
\vect{v}\cdot\vect{v}
=
\|\vect{v}\|^2,
\]

we may also write

\[
\boxed{
\operatorname{proj}_{\vect{v}}\vect{u}
=
\frac{\vect{u}\cdot\vect{v}}
{\|\vect{v}\|^2}
\vect{v}.
}
\]

The result is the component of $\vect{u}$ in the direction
of $\vect{v}$.

\end{frame}


% =========================================================
% ORTHOGONAL DECOMPOSITION
% =========================================================

\begin{frame}{Orthogonal Decomposition}

A vector $\vect{u}$ can be decomposed into two components:

\[
\boxed{
\vect{u}
=
\vect{u}_{\parallel}
+
\vect{u}_{\perp}.
}
\]

Here,
\vspace{-0.1cm}
\[
\vect{u}_{\parallel}
=
\operatorname{proj}_{\vect{v}}\vect{u}
\]

is parallel to $\vect{v}$, while
\vspace{-0.1cm}
\[
\vect{u}_{\perp}
=
\vect{u}-\vect{u}_{\parallel}
\]

is perpendicular to $\vect{v}$.

Therefore,

\[
\boxed{
\vect{u}_{\perp}\cdot\vect{v}=0.
}
\]

This decomposition is fundamental in geometry,
optimization, and machine learning.

\end{frame}


% =========================================================
% DOT PRODUCT AS A LINEAR FUNCTIONAL
% =========================================================

\begin{frame}{Dot Product and Duality}

Fix a vector $\vect{u}$.

The mapping

\[
f_{\vect{u}}(\vect{x})
=
\vect{u}\cdot\vect{x}
\]

takes a vector as input and produces a scalar.

It is linear because

\[
f_{\vect{u}}(\alpha\vect{x}+\beta\vect{y})
=
\alpha f_{\vect{u}}(\vect{x})
+
\beta f_{\vect{u}}(\vect{y}).
\]

Therefore,

\[
\boxed{
\vect{x}\mapsto\vect{u}\cdot\vect{x}
}
\]

is a \textbf{linear functional}.

\end{frame}


% =========================================================
% DUAL SPACE
% =========================================================

\begin{frame}{Dual Space}

The set of all linear functionals on a vector space $V$
forms its \textbf{dual space}, denoted by

\[
\boxed{
V^*.
}
\]

A vector $\vect{u}$ can define a linear functional through

\[
\vect{x}
\mapsto
\vect{u}^{T}\vect{x}.
\]

Thus, the dot product provides a natural connection between
vectors and linear functionals.

\vspace{0.4cm}

\[
\boxed{
\text{Vector}
\quad\longrightarrow\quad
\text{Linear functional}
}
\]

This connection is known as \textbf{duality}.

\end{frame}


% =========================================================
% CROSS PRODUCT
% =========================================================

\begin{frame}{Cross Product}

For vectors
\vspace{-0.1cm}
\[
\vect{u},\vect{v}\in\mathbb{R}^3,
\]

the \textbf{cross product}
\vspace{-0.1cm}
\[
\vect{u}\times\vect{v}
\]

produces another vector in $\mathbb{R}^3$.

If
\vspace{-0.1cm}
\[
\vect{u}
=
\begin{bmatrix}
u_1\\u_2\\u_3
\end{bmatrix},
\qquad
\vect{v}
=
\begin{bmatrix}
v_1\\v_2\\v_3
\end{bmatrix},
\]

then
\vspace{-0.1cm}
\[
\boxed{
\vect{u}\times\vect{v}
=
\begin{bmatrix}
u_2v_3-u_3v_2\\
u_3v_1-u_1v_3\\
u_1v_2-u_2v_1
\end{bmatrix}.
}
\]

\end{frame}


% =========================================================
% CROSS PRODUCT EXAMPLE
% =========================================================

\begin{frame}{Example: Cross Product}

Let
\vspace{-0.1cm}
\[
\vect{u}
=
\begin{bmatrix}
1\\2\\3
\end{bmatrix},
\qquad
\vect{v}
=
\begin{bmatrix}
4\\5\\6
\end{bmatrix}.
\]

Then
\vspace{-0.1cm}
\[
\vect{u}\times\vect{v}
=
\begin{bmatrix}
2(6)-3(5)\\
3(4)-1(6)\\
1(5)-2(4)
\end{bmatrix}.
\]

Therefore,
\vspace{-0.1cm}
\[
\boxed{
\vect{u}\times\vect{v}
=
\begin{bmatrix}
-3\\6\\-3
\end{bmatrix}.
}
\]

\end{frame}


% =========================================================
% CROSS PRODUCT ORTHOGONALITY
% =========================================================

\begin{frame}{Cross Product and Orthogonality}

The cross product is perpendicular to both input vectors:

\[
\boxed{
(\vect{u}\times\vect{v})\cdot\vect{u}=0
}
\]

and

\[
\boxed{
(\vect{u}\times\vect{v})\cdot\vect{v}=0.
}
\]

Therefore,

\[
\boxed{
\vect{u}\times\vect{v}
\perp
\vect{u}
}
\]

and

\[
\boxed{
\vect{u}\times\vect{v}
\perp
\vect{v}.
}
\]

The three vectors therefore define a mutually related
three-dimensional geometry.

\end{frame}


% =========================================================
% MAGNITUDE OF CROSS PRODUCT
% =========================================================

\begin{frame}{Magnitude of the Cross Product}

The magnitude of the cross product is

\[
\boxed{
\|\vect{u}\times\vect{v}\|
=
\|\vect{u}\|
\|\vect{v}\|
\sin\theta.
}
\]

This is equal to the area of the parallelogram formed by
$\vect{u}$ and $\vect{v}$.

\vspace{0.5cm}

Therefore,

\[
\boxed{
\text{Area}
=
\|\vect{u}\times\vect{v}\|.
}
\]

\vspace{0.3cm}

If the vectors are parallel,

\[
\theta=0
\quad\Longrightarrow\quad
\vect{u}\times\vect{v}=\mathbf{0}.
\]

\end{frame}


% =========================================================
% RIGHT-HAND RULE
% =========================================================

\begin{frame}{Direction of the Cross Product}

The direction of

\[
\vect{u}\times\vect{v}
\]

is determined by the \textbf{right-hand rule}.

\vspace{0.4cm}

The cross product is anti-commutative:

\[
\boxed{
\vect{u}\times\vect{v}
=
-(\vect{v}\times\vect{u}).
}
\]

Also,

\[
\boxed{
\vect{u}\times\vect{u}
=
\mathbf{0}.
}
\]

Thus, reversing the order reverses the direction of the
resulting vector.

\end{frame}


% =========================================================
% STANDARD BASIS CROSS PRODUCTS
% =========================================================

\begin{frame}{Cross Products of Standard Basis Vectors}

For the standard basis vectors

\[
\hat{\vect{i}},
\qquad
\hat{\vect{j}},
\qquad
\hat{\vect{k}},
\]

the fundamental relationships are

\[
\boxed{
\hat{\vect{i}}\times\hat{\vect{j}}
=
\hat{\vect{k}}
}
\]

\[
\boxed{
\hat{\vect{j}}\times\hat{\vect{k}}
=
\hat{\vect{i}}
}
\]

\[
\boxed{
\hat{\vect{k}}\times\hat{\vect{i}}
=
\hat{\vect{j}}.
}
\]

Reversing the order changes the sign.

\end{frame}


% =========================================================
% CROSS PRODUCT AS A MATRIX
% =========================================================

\begin{frame}{Cross Product as Matrix Multiplication}

Fix

\[
\vect{u}
=
\begin{bmatrix}
u_1\\
u_2\\
u_3
\end{bmatrix}.
\]

Define the skew-symmetric matrix

\[
[\vect{u}]_{\times}
=
\begin{bmatrix}
0&-u_3&u_2\\
u_3&0&-u_1\\
-u_2&u_1&0
\end{bmatrix}.
\]

Then

\[
\boxed{
\vect{u}\times\vect{v}
=
[\vect{u}]_{\times}\vect{v}.
}
\]

Therefore, for fixed $\vect{u}$, the cross product is a
\textbf{linear transformation} of $\vect{v}$.

\end{frame}


% =========================================================
% LINEAR TRANSFORMATION VIEW
% =========================================================

\begin{frame}{Cross Product as a Linear Transformation}

Fix $\vect{u}\in\mathbb{R}^3$ and define
\vspace{-0.1cm}
\[
T(\vect{v})
=
\vect{u}\times\vect{v}.
\]

Then

\[
T(\alpha\vect{v}+\beta\vect{w})
=
\alpha T(\vect{v})
+
\beta T(\vect{w}).
\]

Therefore,

\[
\boxed{
T:\mathbb{R}^3\rightarrow\mathbb{R}^3
}
\]

is a linear transformation.

\vspace{0.4cm}

Its matrix representation is

\[
\boxed{
T(\vect{v})
=
[\vect{u}]_{\times}\vect{v}.
}
\]

\end{frame}


% =========================================================
% KERNEL OF CROSS PRODUCT TRANSFORMATION
% =========================================================

\begin{frame}{Kernel of the Cross Product Transformation}

Consider

\[
T(\vect{v})
=
\vect{u}\times\vect{v}.
\]

For $\vect{u}\neq\mathbf{0}$,

\[
T(\vect{v})=\mathbf{0}
\]

when $\vect{v}$ is parallel to $\vect{u}$.

Therefore,

\[
\boxed{
\operatorname{Null}(T)
=
\operatorname{span}\{\vect{u}\}.
}
\]

\vspace{0.4cm}

The cross-product transformation therefore collapses
one direction to zero.

This connects directly with the concept of
\textbf{null space}.

\end{frame}


% =========================================================
% RANGE OF CROSS PRODUCT TRANSFORMATION
% =========================================================

\begin{frame}{Range of the Cross Product Transformation}

For
\vspace{-0.1cm}
\[
T(\vect{v})
=
\vect{u}\times\vect{v},
\]

every output is perpendicular to $\vect{u}$.

Therefore,
\vspace{-0.1cm}
\[
\boxed{
\operatorname{Range}(T)
=
\vect{u}^{\perp}.
}
\]

The range is the plane through the origin perpendicular
to $\vect{u}$.

Thus,
\vspace{-0.1cm}
\[
\boxed{
\dim(\operatorname{Null}(T))=1
}
\]

and
\vspace{-0.1cm}
\[
\boxed{
\dim(\operatorname{Range}(T))=2.
}
\]

This satisfies
\vspace{-0.1cm}
\[
1+2=3,
\]

as predicted by the Rank--Nullity Theorem.

\end{frame}


% =========================================================
% DOT VS CROSS PRODUCT
% =========================================================

\begin{frame}{Dot Product vs. Cross Product}

\begin{center}

\begin{tabular}{>{\bfseries}lcc}
\toprule
 & Dot Product & Cross Product\\
\midrule
Input & $\mathbb{R}^n$ & $\mathbb{R}^3$\\
Output & Scalar & Vector\\
Notation & $\vect{u}\cdot\vect{v}$ &
$\vect{u}\times\vect{v}$\\
Geometric role & Angle / projection &
Area / perpendicular direction\\
Zero condition & Orthogonal vectors &
Parallel vectors\\
\bottomrule
\end{tabular}

\end{center}

\vspace{0.4cm}

Both operations provide important geometric information
about vectors.

\end{frame}